\textbf{COLUNI 2023 -- Questão 08}\par\medskip Leia o excerto abaixo da crônica “O futuro na geladeira”, de Moacyr Scliar para responder à questão:

\par

(…) Não foi decisão fácil, como se pode imaginar. Curso de administração ou geladeira? A favor de ambas as coisas, o curso e a geladeira, havia argumentos. O curso era algo com que sonhava havia muito tempo, desde jovem, para dizer a verdade. Primeiro, porque era uma fervorosa admiradora da atividade em si, da administração. Organizar as coisas, fazer com que funcionem, levar uma empresa ao sucesso, mesmo em épocas de crise, sobretudo em épocas de crise, parecia-lhe um objetivo verdadeiramente arrebatador. Com o curso, ela poderia tornar-se, mesmo com idade avançada, numa daquelas dinâmicas executivas cuja foto via em jornais e em revistas. Mas a geladeira... A verdade é que ela precisava de uma geladeira nova. A antiga estava estragada, e tão estragada que o homem do conserto aconselhara-a a esquecer "aquele traste" e partir para algo mais moderno. E isso precisava ser feito com urgência: todos os dias estava jogando fora comida que estragara por causa do inconfiável eletrodoméstico.

\par

O futuro na geladeira. SCLIAR, Moacyr. \textbf{Histórias que os jornais não contam}. 3. ed. Porto Alegre: L\&PM Editores, 2017. p. 23. (Adaptado).

\par

Assinale a alternativa que NÃO apresenta uma compreensão de leitura do fragmento acima:\begin{enumerate}
\item A personagem da crônica se via diante de uma difícil dúvida entre cursar administração ou comprar uma
       geladeira.
\item Existiam argumentos para as duas vontades/necessidades da personagem: a compra da geladeira e fazer
       o curso de administração.
\item Ao cursar administração a personagem poderia realizar um sonho porque poderia contribuir com o
       sucesso de uma empresa no futuro.
\item A questão da geladeira para a personagem poderia esperar porque em função de sua idade avançada o
       curso tornara-se prioridade.\end{enumerate}
